\clearpage
\ExposeTOC
\begin{center}
{\large SGA 7 --- 第 XIX 讲}\par
\vspace{2em}
{\Large\underline{Noether 定理}}\par
\vspace{1.5em}
{\large P. Deligne 著}
\end{center}

\section*{目录}

\noindent 1. 定理的陈述. \hfill 1

\noindent 2. 第 XVIII 讲的一个补充. \hfill 4

\noindent 3. 1.3 的证明. \hfill 5

\noindent 4. 第 XVIII 讲的第二个补充. \hfill 8

\noindent 5. Noether 定理的例外情形. \hfill 11

\section*{1. 定理的陈述}

本讲将 Lefschetz 对下述 Noether 定理 (1.2) 的证明 [3, 第~108 页] 作现代化处理.

\subsubsection*{1.1.}
设 $k$ 为特征 $p$ 的代数闭域. 设 $d\geq1$, 并设
\[
\underline a=(a_i)_{1\leq i\leq d}
\]
为一个各项至少为 $1$ 的整数序列. $k$ 上射影空间 $\mathbb P^{n+d}$ 中维数为 $n$、多重次数为 $\underline a$ 的完全交 $V$, 若在扩张 $k$ 后, 它在射影变换下同构于某个由方程定义的完全交, 而这些方程的系数在素域上代数无关, 则称此完全交为\underline{一般的}（generic）. 以下若无特别声明, 均假设 $d\geq1$, $a_i\geq2$.

\subsubsection*{定理 1.2.}
设 $V\subset\mathbb P^{2+d}$ 为多重次数 $\underline a$ 的一般完全交曲面. 假设
\[
\underline a\ne(2),(3),(2,2).
\]
则 $V$ 上的每条（约化）曲线都是 $V$ 与 $\mathbb P^{2+d}$ 的某个超曲面的概形论交.

说明定理可由其推论
\begin{equation}\tag{1.2.1}
\operatorname{Pic}(\mathbb P^{2+d})=\mathbb Z
\xrightarrow{\ \sim\ }\operatorname{Pic}(V).
\end{equation}
推出. 若 $C$ 为 $V$ 上的曲线, 则 (1.2.1) 蕴含对某个适当的 $n$, 有 $\mathcal O(C)\simeq\mathcal O(n)$. 若 $H$ 为超平面截面, 则
\[
n=\bigl(\textstyle\prod_i a_i\bigr)^{-1}(C,H)>0.
\]
只须证明 $|nH|$ 在 $V$ 上诱导的线性系是完全的. 将 $V$ 写为次数为 $a_i$ 的超曲面 $H_i$ 之交, 置 $V_i=H_1\cap\cdots\cap H_i$. 对 $i$ 归纳可见 $|nH|$ 在 $V_i$ 上诱导的线性系是完全的: 在正合列
\[
H^0(V_i,\mathcal O(n))\longrightarrow
H^0(V_{i+1},\mathcal O(n))\longrightarrow
H^1(V_i,\mathcal O(n-a_{i+1})),
\]
中, 最后一项为零（FAC）.

\subsubsection*{定理 1.3.}
设 $V\subset\mathbb P^{2n+d}$ 为非零偶维数 $2n$、多重次数 $\underline a$ 的一般完全交. 假设满足下列条件之一:
\begin{enumerate}[label=(\alph*)]
\item $k$ 的特征不是 $2$, 且 $(2n;\underline a)$ 不等于 $(2n;2)$, $(2n;2,2)$ 及 $(2;3)$;
\item Hodge 数 $h^{2n,0}$ 非零, 即 $\sum_i(a_i-1)>2n$.
\end{enumerate}
则 $H^{2n}(V,\mathbb Q_\ell(n))$ 中的每个代数上同调类都是线性截面的类 $\eta^n$ 的有理数倍.

由 1.3 推出 (1.2.1). 1.2 中的例外情形恰为 $\sum_i(a_i-1)\leq2$ 的情形; 因此可对 1.2 中的曲面 $V$ 应用 1.3(b).\footnote{原印本在此和式中写作 $d_i$, 虽然多重次数始终记为 $(a_i)$.} 已知 $\operatorname{Pic}^{\tau}(V)=0$（XI 1.8）, 所以 $\operatorname{Pic}(V)$ 单射地映入 $H^2(V,\mathbb Q_\ell(1))$. 由 1.3, $\operatorname{Pic}(V)$ 的每个元素都是 $\mathcal O(1)$ 的类 $\eta$ 的有理数倍, 再由（XI 1.8）得到结论.

显然, 1.2 与 1.3 是否成立并不依赖于所选的一般完全交, 而只依赖于 $p$, $n$ 及 $\underline a$.

\subsubsection*{定理 1.4.}
设 $X\subset\mathbb P$ 为 $k$ 上纯维数 $2n+1$ 的光滑射影簇. 设 $\ell\subset\check{\mathbb P}$ 为 $X$ 的超平面截面的一个 Lefschetz 线束. 若 $k$ 的特征为 $2$, 还假设 $\ell$ 为一般 Lefschetz 线束.

假设 $X$ 满足条件 (LV)（XVIII 5.2.2）, 设 $V$ 为线束 $\ell$ 中的一般超平面截面. 则下列两种情形之一成立:
\begin{enumerate}[label=(\alph*)]
\item $H^{2n}(V,\mathbb Q_\ell(n))$ 的消灭部分 $\operatorname{Ev}(V,\mathbb Q_\ell(n))$ 中, 没有非零上同调类是代数的;
\item 代数上同调类生成 $\operatorname{Ev}(V,\mathbb Q_\ell(n))$.
\end{enumerate}

下面说明 1.4 中“一般”的确切含义.

\begin{enumerate}[label=(\alph*)]
\item 设 $G$ 为对偶于 $\mathbb P$ 的射影空间 $\check{\mathbb P}$ 中直线所成的 Grassmann 簇. $G$ 的“几何点”指从代数闭域的谱到 $G$ 的态射. 每个几何点 $\bar s\to G$ 都定义一个超平面截面线束 $\ell$, 作用于将 $X$ 从 $k$ 标量扩张到 $k(\bar s)$ 所得的簇. 若 $\bar s$ 位于 $G$ 的一般点之上, 则称 $\ell$ 为一般的.
\item 设 $\ell$ 为超平面截面线束（即 $\check{\mathbb P}$ 中的直线）. 每个几何点 $\bar s\to\ell$ 都定义将 $X$ 从 $k$ 标量扩张到 $k(\bar s)$ 所得簇的一个超平面截面. 若 $\bar s$ 位于 $\ell$ 的一般点之上, 则称此截面为一般的.
\end{enumerate}

设 $\bar s$ 如上, 其像为 $\ell$ 的一般点 $s$, 并设 $S\subset\ell$ 为给出奇异超平面截面的有限点集. 单值群 $\pi_1(\ell-S,\bar s)$ 是 $\operatorname{Gal}(k(\bar s)/k(s))$ 的商群, 作用于 $\operatorname{Ev}(V,\mathbb Q_\ell(n))$. 若 $k$ 的特征不是 $2$, 则此作用不可约（XVIII 6.7）; 下文第 2 节将说明, 对一般线束, 此断言不需要特征假设也成立. 由代数上同调类生成的 $\operatorname{Ev}(V,\mathbb Q_\ell(n))$ 的子空间在单值作用（即 Galois 作用）下稳定. 因此它或为零, 或为整个空间, 从而证明 1.4.

\subsubsection*{推论 1.5.}
设 $V\subset\mathbb P^{2n+d}$ 为偶维数 $2n$、多重次数 $\underline a$ 的一般完全交. 则或者
\begin{enumerate}[label=(\alph*)]
\item $H^{2n}(V,\mathbb Q_\ell(n))$ 中的每个代数上同调类都是 $\eta^n$ 的有理数倍; 或者
\item $H^{2n}(V,\mathbb Q_\ell(n))$ 由代数上同调类生成.
\end{enumerate}
若 (b) 成立, 则对特征 $p$ 的代数闭域上的每个维数 $2n$、多重次数 $\underline a$ 的完全交, 群 $H^{2n}(V,\mathbb Q_\ell(n))$ 均由代数类生成.

应用 1.4, 其中取 $X$ 为维数 $2n+1$、多重次数 $(a_1,\ldots,a_{d-1})$ 的一般完全交, 取 $\ell$ 为次数 $a_d$ 的超曲面在 $X$ 上截出的截面的一般线束（参见 XVII 2.4）. 也可取 $X$ 的多重次数为 $(a_1,\ldots,a_d)$, 而 $\ell$ 为超平面截面的一般线束. 最后一个断言由代数上同调类在特殊化后仍为代数类这一事实推出.
